\section{FUTURE WORK}

Although the developed robotic arm successfully performs basic pick-and-place
operations, several improvements can be implemented to enhance system
intelligence, accuracy and industrial applicability.

\subsection{Camera-Based Object Recognition}

A camera-based vision system can replace the current IR-based detection method
to automatically identify the position, shape, and orientation of objects.
This would allow the robotic arm to operate with different object locations
without requiring predefined detection positions.

\subsection{AI-Based Vision System}

Artificial intelligence techniques can be incorporated into the vision system
to enhance the autonomous capabilities of the robotic arm. The proposed
AI-based system can provide the following functionalities:

\begin{itemize}
    \item \text{Object Classification:} Identifying and categorizing
    different types of objects based on their visual characteristics.

    \item \text{Real-Time Tracking:} Continuously tracking the position of
    objects during the robotic operation.

    \item \text{Autonomous Decision Making:} Selecting appropriate actions
    based on the detected object and its position.
\end{itemize}

\subsection{Feedback Control Using Encoders}

Currently, the robotic arm operates using predefined servo positions. Rotary
encoders can be incorporated into the joints to provide real-time position
feedback. This would enable closed-loop control and improve the positioning
accuracy of the robotic arm.

\subsection{Higher Payload Design}

Future improvements can include stronger mechanical structures and higher
torque actuators to increase the payload capacity of the robotic arm. This
would allow the system to handle heavier objects and expand its potential
applications.

\subsection{Wireless Monitoring and Control}

Wireless communication technologies such as Wi-Fi or Bluetooth can be
integrated into the robotic arm to enable remote monitoring and operation.
This would allow users to monitor the system status and control the robotic
arm without requiring a direct wired connection.


